%==============================================================================
%  v9 / sections/intro.tex                                        [owner: A3]
%  Introduction and Related Work are one section as of 2026-07-31; the former
%  sections/related.tex was deleted and its \input removed from the master.
%  Its text is recoverable from git history at the commit before that date.
%==============================================================================

\section{Introduction}\label{sec:intro}

A latent tree graphical model relates $m$ observed leaves through unobserved
internal nodes, and the inference problem is to recover the topology from leaf
data alone \cite{aizenbud2023spectral, semple2003phylogenetics}. Phylogenetics is
the canonical instance and supplies our vocabulary: the leaves are extant
organisms, the internal nodes their extinct ancestors, the data aligned
sequences. Distance-, likelihood- and spectral-based reconstruction are all well
developed \cite{saitou1987neighbor, erdos1999logs, jaffe2021spectral,
mossel2005learning}, but datasets now reach tens of thousands of leaves, and the
accurate methods are the slow ones.

That pressure produced the top-down spectral family, which is spectral clustering
run on a phylogeny. Build a similarity matrix $S$ over the leaves, form the
Laplacian $L=D-S$, split the leaf set by the sign pattern of the Fiedler vector
of $L$, recurse on each side, and merge \cite{aizenbud2023spectral,
vonluxburg2007tutorial, fiedler1975property}. The split is provably consistent:
the two sides are complementary \emph{clans}, the unrooted analogue of a clade
\cite{wilkinson2007clades}, which is what keeps the merge tractable. Both the
error and the cost of the recursion land on the top split, since no parent call
corrects it and it is the only call that sees all $m$ leaves. Every entry of $S$
is a sequence-pair estimate over a length-$\ell$ alignment rather than a lookup,
so that one call costs $\Theta(m^2)$ estimator evaluations --- the cost
divide-and-conquer exists to avoid.

This paper asks what happens to the spectrum when that cost is not paid. Compute
only a random $p$-fraction of the pairs, reweight the survivors by $1/p$, and
take the Fiedler vector of the incomplete matrix, in the manner of low-rank
matrix completion \cite{candes2008exact, candes2009power, candes2009matrix,
recht2011simpler, keshavan2010matrix, bhojanapalli2014universal}. Does its sign
pattern still recover the split, and how small can $p$ be? There is reason to
expect an answer well below $1$: under a molecular clock the population matrix
has exactly two blocks and the Fiedler vector is constant on each clan, so the
signal being sampled is extremely concentrated.

Two things keep this from being routine. What decides recoverability from a
random subset is how spread the leading subspace is across coordinates, not how
large the matrix is \cite{chen2014completing}; and guarantees for recursive
spectral bipartitioning are stated for clusters of comparable size
\cite{Balakrishnan2011clustering, qinrohe2013regularized}, whereas real
phylogenies are markedly more imbalanced than standard null models predict
\cite{blum2006which, yule1925mathematical, kingman1982coalescent,
steel2001properties, aldous2001stochastic}. Imbalance pushes both governing
quantities the wrong way: it concentrates the split signal onto the smaller clan,
raising the sampling cost, and it shrinks the eigenvalue gap that isolates the
Fiedler vector, degrading the signal itself. The two failures are distinct ---
one leaves the eigenvector informative but unaffordable to sample, the other
leaves nothing worth sampling --- so pooling them shows one curve where there are
two mechanisms.

Recovering a split from a sign pattern is a claim about \emph{every} coordinate
of an eigenvector, whereas classical perturbation theory controls only its
Euclidean norm \cite{davis1970rotation};
passing between the two charges every node as though it were the worst one, which
here decides whether a guarantee is satisfiable at all. We therefore work
coordinate-wise \cite{eldridge2017unperturbed, fan2018eigenvector,
abbe2020entrywise, tropp2012userfriendly}.

Our contributions are:
\begin{itemize}
\item We make the imbalance ratio $\eta$ the single parameter through which tree
shape reaches the linear algebra, and express both governing spectral quantities
in it: the subspace coherence in closed form, and a lower bound on the Fiedler
gap (\Cref{sec:bridging}).

\item We prove that a sampling rate of order $\log m/m$ suffices, at fixed
imbalance, for the sign pattern of the empirical Fiedler vector to recover the
split exactly \hp, with an explicit cubic dependence on $\eta$
(\Cref{sec:main}). 

\item The proof works coordinate-wise rather than in Euclidean norm, through a
Neumann expansion and a per-node Bernstein bound (\Cref{sec:outline}); on this
model the Euclidean route costs a further factor of $m$ in the required rate. 

\item The same programme on the double-centred \emph{distance} matrix gives a companion
statement \cite{Griffing2012, BandeltDress1992}, and we test both in simulation,
scoring against the planted split \cite{strehlghosh2002cluster,
fortunato2010community} and stratifying by imbalance rather than pooling.
\end{itemize}

\Cref{sec:problem} fixes the generative model, the observation model and the
two operators we analyse; \Cref{sec:bridging} routes tree shape into the spectrum
through $\eta$. \Cref{sec:main} states the Laplacian rate and \Cref{app:dist} the
distance rate, \Cref{sec:outline} assembles the proof, and \Cref{sec:empirical}
tests them.
